%! Setting Up --- Learning from Data — Unit 8, Lesson 8.0 — Homework (Answer Key)
\documentclass[10pt]{article}
\usepackage{linalg-article}
\usepackage{linalg-key}   % pulls in -boxes; do NOT also load -boxes
\newcommand{\TallMath}[1]{$\displaystyle #1\rule[-1.4em]{0pt}{3.2em}$}
\newcommand{\vv}{\mathbf{v}}
\newcommand{\bb}{\mathbf{b}}
\newcommand{\zero}{\mathbf{0}}
\newcommand{\relu}{\mathrm{ReLU}}

\begin{document}

\pageheader{Unit 8, Lesson 8.0}{Homework --- Answer Key}
\namedateperiod

\vspace{0.12in}

\begin{notesbox}{Practice}
A neural network layer is $\relu(A\vv+\bb)$: a \textbf{linear step} (weight matrix $A$ times the input,
plus a bias $\bb$) then the \textbf{ReLU} $=\max(0,x)$, applied to each entry. Stacking ReLUs builds
\textbf{piecewise-linear} functions; the network \textbf{learns} by shrinking a \textbf{loss} (sum of
squared errors).

\begin{enumerate}[leftmargin=1.9em, itemsep=12pt]
  \item \textbf{Evaluate the ReLU.} Keep positives, replace negatives with $0$.
    \begin{enumerate}[label=(\alph*), leftmargin=1.8em, itemsep=4pt]
      \item $\relu(6)$ \hfill \ans{$6$}
      \item $\relu(-4)$ \hfill \ans{$0$}
      \item $\relu(0)$ \hfill \ans{$0$}
      \item $\relu(2.5)$ \hfill \ans{$2.5$}
      \item $\relu(-1.5)$ \hfill \ans{$0$}
    \end{enumerate}
  \item \textbf{One layer, two inputs.} Use $A=\begin{bmatrix} 1 & -1 \\ 2 & 1 \end{bmatrix}$ and
        $\bb=\begin{bmatrix} 0 \\ -3 \end{bmatrix}$. For each input, compute $A\vv+\bb$, then
        $\relu(A\vv+\bb)$.
    \begin{enumerate}[label=(\alph*), leftmargin=1.8em, itemsep=4pt]
      \item $\vv=\begin{bmatrix} 2 \\ 1 \end{bmatrix}$ \hfill \ans{$A\vv=\begin{bsmallmatrix} 1\\ 5 \end{bsmallmatrix}$, $+\bb=\begin{bsmallmatrix} 1\\ 2 \end{bsmallmatrix}$, $\relu=\begin{bsmallmatrix} 1\\ 2 \end{bsmallmatrix}$ (both positive --- unchanged).}
      \item $\vv=\begin{bmatrix} 1 \\ 3 \end{bmatrix}$ \hfill \ans{$A\vv=\begin{bsmallmatrix} -2\\ 5 \end{bsmallmatrix}$, $+\bb=\begin{bsmallmatrix} -2\\ 2 \end{bsmallmatrix}$, $\relu=\begin{bsmallmatrix} 0\\ 2 \end{bsmallmatrix}$ (ReLU zeroes the $-2$).}
    \end{enumerate}
  \item \textbf{A bending function.} Let $F(x)=\relu(x)-\relu(x-3)$. Evaluate $F(-2)$, $F(1)$, $F(3)$, and
        $F(5)$. Where are the two folds, and describe the shape (what does the graph do after $x=3$)?
        \ansline{$F(-2)=0$, $F(1)=1$, $F(3)=3$, $F(5)=5-2=3$. Folds at $x=0$ and $x=3$. The graph is flat at $0$ until $x=0$, rises with slope $1$ to height $3$ at $x=3$, then stays \emph{flat} at $3$ (a ramp that levels off).}
  \item \textbf{Compare two networks.} The targets are $2,4,5$. Network~A predicts $3,4,6$; Network~B
        predicts $1,4,7$. Compute each loss (sum of squared errors). Which network fits better, and which
        set of weights would gradient descent keep?
        \ansline{Loss$_A=(3{-}2)^2+(4{-}4)^2+(6{-}5)^2=1+0+1=2$. Loss$_B=(1{-}2)^2+(4{-}4)^2+(7{-}5)^2=1+0+4=5$. Network~A fits better (loss $2<5$); gradient descent keeps the lower-loss weights (A).}
  \item \textbf{Explain.} (a) Why is \emph{matrix multiplication} called the ``engine'' of a neural
        network? (b) Why is the ReLU step essential --- what could the network \emph{not} do without it?
        \ansline{(a) Every layer's linear step is $A\vv$ (plus a bias) --- a matrix times a vector (Unit~1) --- so nearly all of a network's arithmetic is matrix multiplication. (b) Without the ReLU the layers are all linear and multiply into one matrix (a single straight-line map), so the network could not bend to fit curved data.}
\end{enumerate}
\end{notesbox}

\begin{extensionbox}
\textbf{Two layers with no ReLU collapse into one (why nonlinearity matters).} Suppose a network used two
linear steps back-to-back with \emph{no} ReLU between them:
$A_1=\begin{bmatrix} 1 & 0 \\ 1 & 1 \end{bmatrix}$ and $A_2=\begin{bmatrix} 1 & 2 \\ 0 & 1 \end{bmatrix}$.
\begin{enumerate}[label=(\alph*), leftmargin=1.8em, itemsep=5pt]
  \item Compute the single matrix $A_2A_1$.
  \item Explain: applying $A_1$ then $A_2$ is the same as applying the one matrix $A_2A_1$ --- still just a
        straight-line map. What must sit \emph{between} the two layers for the network to bend?
\end{enumerate}
\ansline{(a) $A_2A_1=\begin{bsmallmatrix} 3 & 2\\ 1 & 1 \end{bsmallmatrix}$.}
\ansline{(b) A ReLU must sit between them. Its fold makes the middle step nonlinear, so the two layers no longer collapse into a single matrix --- that is what lets the network bend.}
\end{extensionbox}

\begin{spiralbox}
\textbf{Coming next --- Lesson 8.1: Piecewise Linear Learning Functions.} Today one ReLU made one fold.
Next we fold many lines together: a layer $\relu(A\vv+\bb)$ builds a \emph{continuous piecewise-linear}
function, and stacking layers lets it bend enough to fit any data.
\end{spiralbox}

\begin{teachernote}
Item~1 is pure ReLU arithmetic (including a decimal and the $\relu(0)=0$ edge). Item~2 runs a full layer
with a bias on two inputs --- part~(b) makes the ReLU visibly zero a negative entry. Item~3 is the
Lesson~8.1 preview: two folds from two ReLUs make a ramp that levels off. Item~4 uses the loss to
\emph{compare} models (lower loss $=$ better fit, the direction gradient descent moves). Item~5 collects
the two target ideas (matrix multiply is the engine; ReLU supplies the nonlinearity). The extension is the
key ``why nonlinearity'' point: $A_2A_1$ is one matrix, so a ReLU must sit between layers. Most common
slips: applying ReLU to a vector as a whole instead of each entry; forgetting to square (or dropping the
sign before squaring) in the loss; multiplying $A_1A_2$ instead of $A_2A_1$.
\end{teachernote}

\end{document}
